\chapter*{Fiche Objectif}

\noindent Les objectifs principaux du stage sont résumés ci-dessous.

\subsection*{Mise en place d'une plateforme robotique}
Le premier objectif consistait à mettre en place une plateforme de robots holonomes qui servira ensuite au développement d'algorithme d'essaim. Cela a impliqué la migration des robots TurtleBot Mecanum existants de ROS1 vers ROS2, l'adaptation du programme bas niveau avec la bibliothèque Micro-ROS, ainsi que la récupération des donnée du système de capture de mouvement Optitrack/Motive. 
\subsection*{Développement d'un algorithme de contrôle distribué}
Le deuxième objectif portait sur le développement d'un algorithme de contrôle distribué, et de le tester avec la plateforme mise en place. 

\subsection*{Optimisation par commande événementielle}
Le troisième objectif visait à optimiser le contrôle de l'essaim en utilisant des stratégies de commande événementielle. Deux approches ont été implémentées : une commande événementielle classique et une version prédictive.

\subsection*{Analyse de performance et validation}
Le quatrième objectif était dédié à l'analyse des performances et à la validation des algorithmes développés. Des tests ont été menées pour les trois méthodes de contrôle. Les performances ont été évaluées sur la base de la vitesse de convergence, ou encore de l'utilisation des ressources (CPU/RAM). La robustesse du système a également été testée face à des perturbations manuelles.

\subsection*{Préparation du transfert vers des drones}
Enfin, le dernier objectif était de préparer le transfert des travaux vers une plateforme de drones. 


